%=================================%
\section{Discussion}
\label{sec:discussion}
%=================================%


Our methodology provides a robust and self-consistent exploration of the 2D linear torque across various disc parameters and eccentricities. In the previous sections we presented the torque and evolutionary timescale results at some discrete points in the space of disc parameters. In the following Sections \ref{subsection:fitting_functions} and \ref{subsection:migration_tracks} we build upon these results, while providing some additional comparisons and discussion in Sections \ref{subsection:previous_results} and \ref{subsection:caveats}.


%---------------------------------%
\subsection{Fitting functions}
\label{subsection:fitting_functions}
%---------------------------------%

Although the shape of the torque and migration timescale profiles presented in Section \ref{sec:eccentric_torques} appears relatively simple, our attempts to derive simple fitting functions across the three-dimensional $(q,p,h_\mathrm{p})$ parameter space has proven challenging. Indeed, the shape of the $T_\textrm{net}$ curve in Fig.~\ref{fig:torque_eccentricity_h_0.06} is suggestive of a hyperbolic tangent superimposed with a (skewed) Gaussian component. Whilst this ansatz can well describe the curve for a specific panel, there is not a simple relationship linking the fitting parameters across the whole grid as a function of $q$ and $p$ (never-mind the variation associated with the wider data set as a function of $h_\textrm{p}$). This emphasizes the complexity of the general problem and also cautions against the use of simple fitting formulae. 

Instead we have opted to our make data available to community and provide our full data cube of torques and evolutionary timescales in the supplementary material which can be easily read in and interpolated by users for their own purposes. The details of this data are described more in Appendix \ref{app:parameter_space}.

% ---------------------------------- %
\subsection{Migration tracks}
\label{subsection:migration_tracks}
% ---------------------------------- %

To illustrate the usage of our results, we computed some exemplar tracks for the orbital evolution of eccentric planets as a function of time. When doing this, it needs to be remembered that we have computed the torques assuming that the planetary position is held fixed and hence so too are the local disc properties, which in particular scale the dimensional torque through the characteristic flux quantity $F_{J,0}$. However, as the planet moves in $a$, the local aspect ratio will change in general even for a fixed disc profile with constant $(q,p)$ and $h_\mathrm{p}$ corresponding to the initial planetary semi-major axis. To account for this variation of $h$ as the planet migrates, we linearly interpolate our data cube of orbital evolution rates onto the instantaneous values for $(a,e,h(a))$ and the local disc properties. This procedure is detailed more in Appendix \ref{app:parameter_space} where we describe how to access and manipulate our full data cube, with exemplar reading and plotting routines available in the supplementary material. 

With these details in mind, we consider the orbital evolution of an  $M_\mathrm{p}=M_\oplus$ protoplanet around a $M_\odot$ central star. The protoplanet starts at a semi-major axis of $a_0 = 5$ AU with an initial eccentricity $e_0 = 0.12$ (the highest probed in our work) and is embedded in a disc with local surface density $\Sigma(5\textrm{AU}) = 355$ g cm$^{-2}$ and aspect ratio $h(a_0) = 0.062$. Note that we choose this starting aspect ratio to be slightly larger than the minimum value of $h_\textrm{p}=0.06$ so that as the planet migrates inwards, its environment remains within the range of parameters surveyed in Section \ref{sec:eccentric_torques}. This planetary mass is suitably sub-thermal for our adopted aspect ratio, where $M_\mathrm{p}/M_\textrm{th} = 0.0126$. The resulting orbital evolution tracks are plotted in Fig.~\ref{fig:1d_migration} as a function of time, where we see the evolution of $a/a_0$ in the upper panel and $e/e_0$ in the lower panel over 800 orbital timescales $T_0$ (measured with respect to the initial planetary position) or equivalently $9\times 10^3$ yr. Tracks computed for the three different combinations of $(q,p) = (0.0,0.0)$, $(0.0,1.5)$ and $(0.75,0.0)$ are plotted as the black, red and blue lines respectively.
%
%%%%%%%%%%%%%%%%%%%%%%%%%%%%%%%%%
\begin{figure}
    \centering
    \includegraphics[width=\columnwidth]{Figures/1d_migration.pdf}
    \caption{Migration tracks for an Earth mass planet embedded in a disc at initially $a_0 = 5$AU, where $\Sigma(5\textrm{AU}) = 355$ g cm$^{-2}$ and $e_0 = 0.12$. \textit{Upper panel}: semi-major axis evolution. \textit{Lower panel}: eccentricity evolution. Time is measured in orbital periods $T_0$ at $a_0$. The black, red and blue curves denote disc profiles which adopt $(q,p) = (0.0,0.0)$, $(0.0,1.5)$ and $(0.75,0.0)$ respectively.}
    \label{fig:1d_migration}
\end{figure}
%%%%%%%%%%%%%%%%%%%%%%%%%%%%%%%%%

We see that over the explored time interval, the semi-major axis changes very little whilst the eccentricity varies quite dramatically. Orbital elements vary the most at the very beginning of the tracks, as the planet progresses through a range of eccentricities. In the uniform density discs (black and blue curves), we observe a brief spell where the planet migrates outwards in accordance with the semi-major axis drift rate reversals seen in panels (a) and (m) of Fig.~\ref{fig:torque_eccentricity_h_0.06}. However, this trend quickly reverses as the eccentricities decrease and the typical inwards migration resumes -- in fact peaking around the transonic value of $e$. The three curves split into distinct paths as the varying disc properties separate the migration rates, with the steepest surface density discs leading to the quickest inwards drift. Meanwhile, the eccentricity paths almost overlap suggesting that the eccentricity damping is only weakly sensitive to the details of the disc profile, as suggested by the quantitatively similar $\tau_e^{-1}$ curves seen across all panels in Fig.~\ref{fig:eccentric_timescales_h_0.06}. 

%%%%%%%%%%%%%%%%%%%%%%%%%%%%%%%%%
\begin{figure}
    \centering
    \includegraphics[width=\columnwidth]{Figures/a_e_tracks.pdf}
    \caption{Migration tracks in $(a,e)$ phase space for planets in a uniform surface density, globally isothermal disc with $a_0 = 5$AU and $\Sigma(5\mathrm{AU}) = 355$ g cm$^{-2}$. The different tracks correspond to different initial eccentricities $e_0$ and the colour of each curve indicates the time in units of the orbital period at $a_0$. The dashed vertical line denotes the initial semi-major axis of the planet. The inset shows the time it takes for each track to reach $a=0.995a_0$ as a function of $e_0$.}
    \label{fig:a_e_tracks}
\end{figure}
%%%%%%%%%%%%%%%%%%%%%%%%%%%%%%%%%

Another way to visualise this evolution is to consider tracks in $(a,e)$ space as shown in Fig.~\ref{fig:a_e_tracks}. Here, we once again assume the same system properties as for the black curve in Fig.~\ref{fig:1d_migration} wherein the disc has a uniform temperature and surface density profile. However, now the $M_\mathrm{p}=M_\oplus$ planet is initialised at $5$ AU (along the vertical dashed grey line) with different eccentricities ranging from $e_0=0.02$ to $0.12$ in steps of $0.02$. The migration is followed for each of these initial conditions for 350 orbital timescales at the starting semi-major axis ($\sim 3.9\times 10^3$ yr). The time along each track is visualised by the colour. Notably, for the most eccentric initial condition we see a brief period of outwards migration with continued eccentricity damping. This trend rapidly reverses once the planetary eccentricity has diminished, and the planet then starts migrating inwards, closely overlying the curve for the $e=0.10$. 

The initial eccentricity clearly segregates the time at which the different curves reach a fixed value of $a$. For example, the inset panel shows the time at which tracks of varying $e_0$ cross $a/a_0 = 0.995$. We see that this time reaches a minimum for starting eccentricities $e_0$ are around the peak of $\tau_a^{-1}$, as seen in Fig.~\ref{fig:eccentric_timescales_h_0.06}. Meanwhile, it takes longer to reach $0.995a_0$ for the tracks with lower and higher $e_0$, for which the migration in less efficient (or even reverses sign).

%---------------------------------%
\subsection{Comparison with previous prescriptions}
\label{subsection:previous_results}
%---------------------------------%

We now provide a detailed comparison of our eccentric planet-disc coupling results with past studies. As pointed out by \cite{IdaEtAl_2020}, the existing literature is complicated by a host of studies which probe different disc regimes and employ various methodologies and approximations. Often these approaches are stitched together in a heuristic fashion to motivate torque fitting formulae over a wider parameter space. Inevitably such a procedure has led to a variety of possible planet migration characteristics depending on which model is adopted. For these reasons, here we will focus on the comparison with two previous studies which are most relatable to our current approach -- namely they investigate 2D discs and consider torque behaviour across a range of eccentricities, in particular pushing towards supersonic values. 

Firstly, we summarise the approach of \cite{PapaloizouLarwood_2000} (hereafter PL00) who considered discs with $q=1.0$, $p=1.5$ and aspect ratio $h_\textrm{p} = 0.05$ (uniform across the disc for the adopted $q$). Note, that this choice of surface density profile is often targeted by previous studies in order to suppress the complicating effects of corotation resonances (see Section \ref{subsubsec:corotation}). However, although this choice of $p$ removes the leading order corotation term associated with vortensity gradients, it does not remove the entropy gradient contributions mentioned in Section \ref{subsection:torque_trends}. Whilst they also extract the full $(m,l)$ Fourier decomposition of the perturbing potential, instead of solving for the wave mode structure numerically they exploit the modified WKB formula of \cite{Artymowicz_1993}, as presented earlier in equation \eqref{eq:L_WKB}. They provide fits for the resulting angular momentum and eccentricity damping rates to be
%
\begin{equation}
    \label{eq:tau_m_inv_PL00}
    \tau_{L,\textrm{PL00}}^{-1} = 7.33 \left(\frac{b}{0.5}\right)^{-1.75}h_\textrm{p}\frac{1-(\Tilde{e}/1.1)^4}{1+(\Tilde{e}/1.3)^5}\tau_0^{-1} ,
\end{equation} 
%
\begin{equation}
    \label{eq:tau_e_PL00}
    \tau_{e,\textrm{PL00}}^{-1} = 4.26 \left(\frac{b}{0.5}\right)^{-2.5}\left(1+\frac{1}{4}\tilde{e}^3\right)^{-1}h_{\textrm{p}}^{-1}\tau_0^{-1},
\end{equation}
%
as a function of softening  parameter $b$, aspect ratio $h_\mathrm{p}$ and, crucially, the scaled eccentricity $\tilde{e}$. This early work clearly predicts a torque reversal when $\tilde{e}\sim 1.1$, in keeping with our findings also. Furthermore, the eccentricity damping rate decreases $\propto \tilde{e}^3$ for high eccentricities. The PL00 orbital evolution rates (\ref{eq:tau_m_inv_PL00})-(\ref{eq:tau_e_PL00}) are illustrated in Fig.~\ref{fig:tau_m_comparison} with green dashed curves.

Secondly, \cite{CresswellNelson_2006} (hereafter CN06) performed 2D, hydrodynamic simulations for a range of planetary eccentricities also embedded in discs with $q=1.0$, $p=1.5$ and $h_\textrm{p} = 0.05$. They directly measured the angular momentum and eccentricity damping rates which were compared to the predictions of PL00. They state that their measured angular momentum damping rates at low $e$ are a factor of $\sim 3$ slower than the PL00 prediction, whilst the high $e$ rates are a factor of $\sim 3/2$ slower. Meanwhile, they find good agreement with the predicted eccentricity damping timescale for high $e\gtrsim 0.08$ but report a clear discrepancy for $e\lesssim 0.05$ where PL00 again give a rate roughly 3 times faster. CN06 also emphasize the sensitivity to the softening parameter and find that a value of $b\sim 0.5-0.6$ gives good agreement with the 3D analytical results of \cite{TanakaEtAl_2002} for the migration rate. 

% It should be noted that CN06 are not entirely consistent in this comparison since they convert the PL00 angular momentum damping formula into a semi-major axis migration rate via $\tau_a^{-1} = 2 \tau_{L}^{-1}$, which is only true for $e=0.0$ (see equation \eqref{eq:timescale_relation}). %

Informed by these simulations, they combine in a somewhat ad-hoc fashion 3D analytical calculations presented by \cite{TanakaEtAl_2002} and \citet{TanakaWard_2004} with the eccentricity dependencies found by PL00. Their suggested angular momentum damping rate prescription is 
%
\begin{equation}
    \label{eq:tau_m_inv_cn}
    \tau_{L,\textrm{CN06}}^{-1} = \frac{2.7+1.1 p}{2} h_\textrm{p} \frac{1-(\Tilde{e}/1.1)^4}{1+(\Tilde{e}/1.3)^5}\tau_0^{-1} ,
\end{equation}
%
which apparently gives very good agreement with their 2D simulation measurements for low $e$ and decent agreement at large $e$. Meanwhile, their final eccentricity prescription is more questionable and given by
%
\begin{equation}
\label{eq:tau_e_inv_cn}
    \tau_{e,\textrm{CN06}}^{-1} = 
    \frac{0.78}{Q_e} \left(1+\frac{1}{4}\tilde{e}^3\right)^{-1} h_{\textrm{p}}^{-1}\tau_0^{-1} ,
\end{equation}
%
where $Q_{e}\sim 0.1$ was chosen by \cite{CresswellNelson_2006} to approximately model their simulation results. The CN06 orbital evolution rates (\ref{eq:tau_m_inv_cn})-(\ref{eq:tau_e_inv_cn}) are also shown in Fig.~\ref{fig:tau_m_comparison}, for different $Q_e$ in panel (b).

Given that CN06 visually demonstrates good agreement between the simulations and PL00 for high $e$, we find that the choice of $Q_e \sim 0.2$ brings equations \eqref{eq:tau_e_PL00} and \eqref{eq:tau_e_inv_cn} into better accordance, see blue dashed curve in Fig.~\ref{fig:tau_m_comparison} (b). Such a prescription still disagrees significantly with the low $e$ simulation results for $\tau_{e}^{-1}$, as per the comparison with the PL00 prescription discussed previously. Therefore CN06 implement a constant rate below $e<0.08$ to limit the damping rate. In fact, if one instead adopts $Q_e = 1$ (orange dashed curve in Fig.~\ref{fig:tau_m_comparison}(b)) then the low eccentricity limit of equation \eqref{eq:tau_e_inv_cn} becomes  equivalent to the analytical result proposed by \cite{TanakaWard_2004} for small eccentricity damping rates in a 3D disc. It might be hoped that a low-$e$, 2D setup with $b \sim 0.5$ would then agree better with this analytical result, akin to the aforementioned match found for $\tau_{L}^{-1}$ between the softened 2D simulations of CN06 and the 3D migration rates presented by \cite{TanakaEtAl_2002}.

%%%%%%%%%%%%%%%%%%%%%%%%%%%%%%%%%
\begin{figure}
    \centering
    \includegraphics[width=\columnwidth]{Figures/tau_comparison.pdf}
    \caption{Migration rates for a disc with $q = 1.0$, $p = 1.5$, $h_\textrm{p} = 0.05$ and $b=0.5$. Our data is plotted as the black points and connecting points. \textit{Upper panel (a)}: Comparison of angular momentum damping rates $\tau_{L}^{-1}$. Green dashed curve is the PL00 prediction given by the equation \eqref{eq:tau_m_inv_PL00}. Blue dashed curve is the CN06 prescription given by the equation \eqref{eq:tau_m_inv_cn}. \textit{Lower panel (b)}: Comparison of eccentricity damping rates $\tau_e^{-1}$. Green dashed curve is for PL00 prescription, see equation \eqref{eq:tau_e_PL00}. The red, blue and orange dashed lines plot the CN06 prescription given by the equation \eqref{eq:tau_e_inv_cn} for $Q_e =$ 0.1, 0.2 and 1.0 respectively.}
    \label{fig:tau_m_comparison}
\end{figure}
%%%%%%%%%%%%%%%%%%%%%%%%%%%%%%%%%


To compare our methodology with these previous studies, we use our pipeline to extract the torque for a comparable disc with $q=1.0$, $p=1.5$, $h_\textrm{p} = 0.05$ and $b = 0.5$. We then compute $\tau_{L}^{-1}$ and $\tau_e^{-1}$ as a function of planetary eccentricity and compare them with the literature prescriptions in Fig.~\ref{fig:tau_m_comparison}. We plot our results as the black data points connected by a solid line. In panel (a), we compare with the PL00 prediction \eqref{eq:tau_m_inv_PL00} and the CN06 prescription \eqref{eq:tau_m_inv_cn}. The CN06 result seems to qualitatively agree with our results, especially at low $e$, although some discrepancies are obvious in the transonic and high-$e$ regimes. By extension, our results should also qualitatively agree well with the 2D simulations of CN06. On the other hand, the PL00 curve deviates from our calculation by a factor of $\sim 2.5$ in the low-$e$ limit and $\sim 1.7$ in the high-$e$ limit. This discrepancy is reminiscent of similar differences found between PL00 and the simulations of CN06. Hence our method aligns better with the numerics and justifies the use of our global linear theory. Namely, it emphasizes the importance of calculating the full global structure of the excited wave modes if one cares about order unity differences, rather than simply the localised WKB approach.

In the lower panel (b), we instead compare the eccentricity damping rates. Our results are once again plotted as black dots connected by the solid line, which clearly exhibits a characteristic enhancement around the transonic value of $e$. This enhancement in the eccentricity damping timescales (as well as semi-major axis damping in Fig.~\ref{fig:eccentric_timescales_h_0.06}), is robustly found in simulations of eccentric planets embedded in discs \citep[e.g. CN06,][]{Cresswell2007}. However, this effect is typically not captured by orbital migration fitting functions presented elsewhere in the literature \citep{IdaEtAl_2020}. For example, the PL00 prescription (equation \eqref{eq:tau_e_PL00}) is plotted as the dashed green line. In keeping with the numerical comparison of CN06, this gives very good agreement for supersonic values $e$. However, for low values of $e$ the PL00 result gives a damping rate, which is faster than ours by a factor of $\sim 4$. This echoes the discrepancy found by CN06 who reported the eccentricity damping in their simulations to be a factor of $\sim 3$ slower than PL00 and thus in much better alignment with our results. The CN06 model given by equation \eqref{eq:tau_e_inv_cn} is also over-plotted for three different values of $Q_{e} = (0.1,0.2,1.0)$ as the red, blue and orange lines respectively. The $Q_e = 0.2$ result give close correspondence with the PL00 result and shows the expected good agreement with our results at high $e$. Meanwhile, the $Q_e = 1.0$ result gives better agreement for the lower eccentricity end. Our result, which is in closer agreement with the simulations, seems to be an interpolation of these two limiting cases (i.e. may be modelled using $e$-dependent $Q_e$, dropping as $e$ increases). It should be emphasized that whilst these previous prescriptions perform well in limited eccentricity intervals, only our self-consistent result smoothly connects across the transonic regime and captures both qualitative and quantitative consistency with past simulations.

Finally, it should be noted that our method also yields a semi-major axis damping rate which remains positive for the adopted $p$ and all examined eccentricities and therefore does not exhibit outwards migration. Although not plotted here, we find that $\tau_a^{-1}$ is levelling off towards a positive value beyond $e>0.12$, which is consistent with the fact that CN06 report no outwards migration of the planet even for large values of eccentricity. This is in stark contrast with the results of PL00 which predicted expansion of the planetary orbit for transonic eccentricities (as seen in Fig. 1 of \cite{IdaEtAl_2020}).

%---------------------------------%
\subsection{Caveats of our approach}
\label{subsection:caveats}
%---------------------------------%

Our approach provides a self-consistent calculation of the linear torque across a wide disc parameter space and range of eccentricities. However, we acknowledge that this scope comes at the expense of a number of compromises which we now mention.

Firstly, we should remember that we are calculating the response for a 2D disc. The process of vertical averaging motivates a softened potential, which introduces the somewhat arbitrary $b$ parameter. This is typically tuned by comparison of 2D results with 3D simulations \citep[e.g.][]{CresswellNelson_2006}. Notably, \cite{Cresswell2007} and \cite{DAngeloLubow_2010} performed both 2D and 3D simulations showing that the 2D torque can be brought into qualitative agreement with the 3D value for specific choices of $b$.  However, the variation of torque with disc characteristics $p$ and $q$ is also sensitive to the exact value of $b$ chosen. More recent numerical simulations of \cite{JimenezMasset_2017a} corroborate these results in the locally isothermal case.

There has also been analytical work extending the planet-disc interaction to three dimensions. In particular, \cite{TanakaEtAl_2002} solved for the 3D wave modes and calculated the resulting migration rate for a circular planet, avoiding the need for a softening parameter. Recently this has been extended to the locally isothermal case \citep{TanakaOkada_2024} which compares favourably with the simulations discussed above. However, although they considered an eccentric perturber with small $e$ in \cite{TanakaWard_2004}, this did not allow them to probe the supersonic regime of key interest in this work.

Another notorious complication is the treatment of the corotation torque. We have limited our attention to the case of the linear corotation torque. However, the evanescent waveform launched in the vicinity of corotation resonances \citep{GoldreichTremaine_1979} means that the angular momentum is deposited locally and might modify the background disc structure. This develops as a nonlinear \textit{horseshoe} drag \citep[e.g.][]{Ward_1991, PaardekooperEtAl_2010} wherein fluid elements exchange angular momentum with the planet as they undergo librating orbits in the vicinity of corotation. If there is sufficient viscosity in the disc to communicate away the deposited angular momentum flux, it is reasoned that the linear corotation result can still be a good approximation. Indeed, \cite{OgilvieLubow_2003} showed that the classic linear result can be recovered for non-coorbital eccentric corotation resonances when the perturbing potential strength relative to the dissipation is small. We can therefore speculate that as the eccentricity of the perturber increases and the the non-coorbital resonances become more important, the linear corotation torques may dominate over the nonlinear horseshoe contributions. 

Whilst these caveats will change some of the quantitative details for 3D discs which permit corotation nonlinearities, the qualitative behaviour we find seems to be borne out in previous simulations. We encourage comparison with targeted future simulations of 2D and 3D locally isothermal eccentric planet-disc interactions to further test our results.

% \cite{PaardekooperEtAl_2010} - linear calculations of circular planet with locally isothermal profile. Get a torque formula from 2D softened wave approach. Not clear how the corottaion is being decomposed 

% \cite{TanakaOkada_2024} - 3D locally isothermal calculation with Landau prescription to deal with irregular singularity. Builds on their previous work.

% \cite{DAngeloLubow_2010} - 3D simulations of locally isothermal disc. Compared with softened locally isothermal predictions of Menou and Goodman suggests discrepancy. 
% They perform 2D simulations as well and there is a comment that the coefficients of the torque are sensitive on the softening length as well. Can get qualittaive agreement but often the details of the torque dependence on disc properties can be quantitatively different between 2D and 3D.

% \cite{MenouGoodman} - presented a softened 2D locally isothermal disc torque. Only present the Lindblad component and really. But this includes the effect of surface temperature only in a local fashion but plugged into the WKB modified formulae of Ward. This is a bit like a locally isothermal version of the method used by Papaloizou and Larwood.







